\documentclass[a4paper, 12pt]{article}
\usepackage[slovene]{babel}
\usepackage{lmodern}
\usepackage[T1]{fontenc}
\usepackage[utf8]{inputenc}
\usepackage{url}
\usepackage{xcolor}
\usepackage[square,numbers,sort&compress]{natbib}
\bibliographystyle{elsarticle-num}

\definecolor{munsell}{rgb}{0.0, 0.5, 0.69}
\newcommand\cmnt[1]{\textcolor{munsell}{#1}}


\topmargin=0cm
\topskip=0cm
\textheight=25cm
\headheight=0cm
\headsep=0cm
\oddsidemargin=0cm
\evensidemargin=0cm
\textwidth=16cm
\parindent=0cm
\parskip=12pt

\renewcommand{\baselinestretch}{1.2}

\begin{document}

%%%%%%%%%%%%%%%%%%%%%%%%%% Izpolni kandidat! %%%%%%%%%%%%%%%%%%%%%%%%%%
\newcommand{\ImeKandidata}{Aleksander} % Ime
\newcommand{\PriimekKandidata}{Kovač} % Priimek
\newcommand{\VpisnaStevilka}{63220478} % vpisna številka
\newcommand{\StudijskiProgram}{Multimedija, MAG} % Študijski program/smer
\newcommand{\NaslovBivalisca}{} % kaniddatov naslov
\newcommand{\SLONaslov}{Sistem za vadbo biljard igre osmica v obogateni resničnosti} % naslov dela v slovenščini
\newcommand{\ENGNaslov}{Augmented Reality Training System for Eight-Ball Pool} % naslov dela v angleščini
%%%%%%%%%%%%%%%%%%%%%%%%%% Konec izpolnjevanja %%%%%%%%%%%%%%%%%%%%%%%%%%


\newcommand{\ProblemParagraph}{\noindent\textbf{Problem in pregled stanja tehnike.} }
\newcommand{\ContributionsParagraph}{\noindent\textbf{Pričakovani znanstveni ali tehnični prispevki. } }
\newcommand{\MethodologyParagraph}{\noindent\textbf{Metodologija in ovrednotenje. } }

\newcommand{\Kandidat}{\ImeKandidata~\PriimekKandidata}
\noindent
\Kandidat\\
\NaslovBivalisca \\
Študijski program: \StudijskiProgram \\
Vpisna številka: \VpisnaStevilka
\bigskip

{\bf Komisija za študijske zadeve}\\
Univerza v Ljubljani, Fakulteta za računalništvo in informatiko\\
Večna pot 113, 1000 Ljubljana\\

{\Large\bf
{\centering
    Vloga za prijavo teme magistrskega dela \\%[2mm]
\large Kandidat: \Kandidat \\[10mm]}}


\Kandidat, študent/-ka magistrskega programa na Fakulteti za računalništvo in informatiko, zaprošam Komisijo za študijske zadeve, da odobri predloženo temo magistrskega dela z naslovom:

%\hfill\begin{minipage}{\dimexpr\textwidth-2cm}
Slovenski: {\bf \SLONaslov}\\
Angleški: {\bf \ENGNaslov}
%\end{minipage}

Tema je bila že potrjena lani in je ponovno vložena: {\bf \textit{DA} }

Izjavljam, da so spodaj navedeni mentorji predlog teme pregledali in odobrili ter da se z oddajo predloga strinjajo.

Magistrsko delo nameravam pisati v slovenščini. % In case you would like to write the thesis in English, comment this line out, and use the following template to explain your request:
%Komisijo zaprošam, da odobri pisanje magistrskega dela v angleškem jeziku z obrazložitvijo ... .

Za mentorja/mentorico predlagam:

%%%%%%%%%%%%%%%%%%%%%%%%%% Izpolni kandidat! %%%%%%%%%%%%%%%%%%%%%%%%%%
\hfill\begin{minipage}{\dimexpr\textwidth-2cm}
Ime in priimek: izr. prof. dr. Matevž Pesek\\
Ustanova: Fakulteta za računalništvo in informatiko, Univerza v Ljubljani\\
Elektronski naslov: matevz.pesek@fri.uni-lj.si
\end{minipage}
%%%%%%%%%%%%%%%%%%%%%%%%%% Konec izpolnjevanja %%%%%%%%%%%%%%%%%%%%%%%%%%

\bigskip

\hfill V Ljubljani, \today.\newpage

\section{Ključne besede}

slovensko: obogatena resničnost, zaznava predmetov, sledenje predmetom, namizna igra \\
angleško: augmented reality, object detection, object tracking, tabletop game

\section*{Podroben predlog magistrske naloge}
\ProblemParagraph~Trenutno prevladujoče rešitve za podporo učenju biljarda temeljijo na dvorazsežnostnih nadglavnih kamerah ali globinskih sistemih ter tudi naglavnimi prikazovalniki \cite{Sousa2016,WenKai2024,Yan2024Enhancing, BWisePool2025}. Čeprav izboljšajo uspešnost začetnikov, podpirajo le enega uporabnika, zahtevajo drago strojno opremo in ne zagotavljajo stabilne poravnave med navideznimi in realnimi objekti, kar vodi do napak pri udarcih. Metode, uporabljene v projektu »pix2pockets« \cite{schiøtt2025pix2pocketsshotsuggestions8ball}, dosegajo zaznavo z napako do 4 milimetre, vendar ne naslovijo prikaza v mešani resničnosti in vendar te rešitve ne ponujajo vpogleda v to, kako bi se takšna natančnost ohranila pri sočasni uporabi več naglavnih prikazovalnikov ali pri pravilnem zastiranju med navideznimi in realnimi objekti. Zato ostaja odprto vprašanje, ali je mogoče s cenovno dostopnimi naglavnimi prikazovalniki (angl. ``head-mounted display''), snemalnimi napravami (npr. pametni telefon) in (prenosnimi) osebnimi računalniki vzpostaviti večuporabniški sistem, ki zanesljivo prikaže navidezne objekte pomešane z resničnimi, na primeru igre biljard.

\ContributionsParagraph~Namen naloge je razvoj prototipa za vadbo igre osmica v mešani resničnosti, ki bo omogočal hkratni prikaz istega navideznega stanja mize na dveh naglavnih prikazovalnikih Meta Quest~3. Prispevek vključuje postopek za preslikavo  dvorazsežnostnih zaznav krogel in žepov v tridimenzionalni koordinatni sistem z jasno definiranimi metrikami napake ter uvedbo osnovnega postopka zastiranja v načinu transparentnega pogleda (angl. ``passthrough''). \cite{Macedo2023Occlusion,Krajancich2020Factored,Zhang2023AddOn,MetaDevelopers2024b}. Dodatno bomo pripravili nabor enostavnih ročnih kretenj kot osnovo za kasnejše sledenje telesu noter-navzen (angl. ``inside-out body tracking'') \cite{UploadVR2023} ter protokol za ovrednotenje učnih učinkov v kontekstu izobraževalnih in športnih aplikacij \cite{Chow2023Exergaming}.

\MethodologyParagraph~Metodologija bo obsegala umerjen zajem nad standardno mizo za biljard, implementacijo dvorazsežnostnega algoritma zaznave krogel \cite{schiøtt2025pix2pocketsshotsuggestions8ball}, preslikavo rezultatov v trorazsežnostni prostor ter validacijo začetnih in končnih položajev z referenčnimi meritvami. V ogrodju razvitega s strani podjetja Meta, bomo implementirali prikaz navideznih elementov v načinu tranparentenga pogleda, osnovni mehanizem zastiranja in preproste ročne kretnje. Večuporabniški način bomo vzpostavili preko povezave z računalnikom, ki bo na primarni prikazovalnik pošiljal lokacije krogel, sekundarni prikazovalnik pa bodo prejemal informacijo preko primarnega. Sistem bomo ovrednotili s kvantitativnimi metrikami poravnave in kvalitativnim testiranjem udeležencev različnih sposobnostnih stopenj.
{\small
\bibliography{bib/references}
}

\end{document}


